\newpage

\ifdefined\useChapters
\chapter{Reflexo}
\else
\chapter{}
\fi

\begin{center}
\textit{Helena e o garçom}
\end{center}

\textbf{O garçom}

O restaurante amanheceu silencioso. O La Fontaine cheirava ao vinho derramado na noite anterior e ao produto barato que o garçom usava para limpar o mármore. Ele já havia passado a toalha três vezes na mesa do canto, embora a mancha tivesse saído na primeira. Não queria admitir que continuava ali por causa da taça de Helena.

Na mesa do canto, a taça que Helena havia usado ainda estava ali.  
Ele a pegou com cuidado, girando o vidro entre os dedos.  
A luz atravessou o cristal e desenhou uma faixa torta no balcão. Nada aconteceu. Mesmo assim, o coração dele, velho e ranzinza, bateu mais forte do que devia. Colocou a taça contra a luz outra vez, irritado consigo mesmo por esperar algum sinal.

\begin{center}
* * *
\end{center}

\textbf{Helena} acordou tarde.  
O quarto estava escuro, as cortinas cerradas, e o relógio digital marcava onze e vinte e três. Helena o conferiu duas vezes, como se a hora comum pudesse garantir que a noite anterior não havia acontecido.
Passou a mão no rosto e sentiu a pele diferente, mais lisa, mais viva.  
O espelho do banheiro a esperava como um velho inimigo.

Encostou-se na pia e ficou ali, observando a si mesma.  
As olheiras, os ombros pesados, o olhar cansado — e, ainda assim, algo novo.  
No fundo dos olhos, havia uma faísca, um lampejo de possibilidade.

Fechou os olhos e pensou na noite passada, na cena impossível, no garoto que sumira.  
O rosto que vira no vidro não era o seu, mas havia obedecido. Se conseguira uma vez, conseguiria de novo. Trancou a porta do banheiro antes de tentar.

\begin{center}
* * *
\end{center}

\textbf{O garçom} atendeu um cliente habitual, um homem que entrava todas as tardes, sentava na mesma cadeira e pedia o mesmo vinho antes de abrir o cardápio.

— O de sempre — disse o homem, colocando o telefone sobre a mesa.

O garçom segurou a garrafa, mas não serviu.

— Tem certeza? O senhor reclamou dele na semana passada.

Era mentira. O cliente olhou a garrafa, depois o cardápio, depois o próprio telefone, como se algum dos três pudesse corrigir sua memória.

— Eu reclamei?

— Bastante.

O homem desistiu do vinho e pediu água. Quando foi pagar, errou duas vezes a senha que usava todos os dias. Na terceira tentativa, perguntou ao garçom se aquela era mesmo a máquina do restaurante.

O garçom não entendia bem o que acontecia, mas reconheceu um prazer silencioso crescendo dentro dele.
O prazer não vinha apenas de confundir o cliente. Vinha de vê-lo desistir da própria lembrança e entregar ao homem de uniforme a autoridade de dizer o que havia acontecido. Bastaram duas frases e a certeza do homem começou a cair aos pedaços. O garçom aproximou a máquina mais uma vez, só para ver até onde aquilo ia.

Ele deu um passo para trás, ofegante.  
O cliente foi embora sem conferir o troco. O garçom contou as moedas devagar. Havia devolvido uma a menos de propósito e, ainda assim, o homem agradecera.

Ficou parado, o coração batendo forte demais.  
Olhou o cliente e sussurrou:  
— O que é isso...?

Mas, dentro dele, a pergunta era outra: por que passara tantos anos sendo ignorado, se podia fazer qualquer um desconfiar até do próprio nome? E por que aquela descoberta parecia mais justa do que deveria?

\begin{center}
* * *
\end{center}

\textbf{Helena} se sentou na frente do espelho e começou o experimento.  
Respirou fundo, ajeitou o cabelo.  
— Eu estou diferente — disse em voz baixa.  
A frase soou estranha no ar, mas o espelho respondeu.  
Por uma fração de segundo, seu rosto afinou, os olhos clarearam, os lábios se curvaram num sorriso que não era dela.  

O reflexo a olhou de volta, vivo.  
Assustada, ela piscou — e tudo voltou ao normal.  
Encostou a mão no vidro; o espelho estava quente.  
Sorriu.  
— Então é assim...

Passou o resto da manhã testando. Afinou o nariz, desfez. Clareou os olhos, perdeu a concentração quando o telefone tocou. Tentou mudar o cabelo inteiro e conseguiu apenas uma mecha loira que desapareceu assim que riu. Anotou os resultados no verso de uma conta: precisava dizer em voz alta, precisava visualizar cada detalhe e, principalmente, não podia hesitar. Depois da quinta tentativa, as pernas tremiam como se tivesse subido dez lances de escada, mas ela continuou.

\begin{center}
* * *
\end{center}

Enquanto isso, o \textbf{garçom} foi até a porta do restaurante com um bloco de comandas no bolso. Uma mulher discutia com o filho porque ele queria atravessar a rua sozinho. O garçom se aproximou e perguntou se ela tinha certeza de que o sinal estava aberto. A mulher olhou para a luz verde, puxou o menino para trás e esperou o próximo ciclo. Um homem ao telefone dizia que havia estacionado naquela rua; depois de cruzar o olhar do garçom, apertou o alarme de quatro carros diferentes.

Gostou da sensação.  
Não era poder de criar.  
Era o poder de fazer alguém parar diante de uma porta aberta. De olhar a própria assinatura e não reconhecê-la. De trocar uma certeza por um “talvez” e deixar que a pessoa fizesse o resto sozinha.

Encostou a mão na vidraça e encontrou o próprio reflexo: uniforme manchado, cabelo ralo, boca apertada. Durante anos, aquele rosto só aparecia quando alguém queria reclamar da conta. Agora duas pessoas estavam paradas na calçada por causa de uma pergunta dele. Sorriu, e dessa vez gostou do homem que sorriu de volta.

\begin{center}
* * *
\end{center}

\textbf{Helena} saiu para caminhar à tarde.  
O sol parecia mais claro do que nunca, e o vento morno a fazia sentir-se leve.  
Testou uma ideia: imaginou que as pessoas a reconheceriam como alguém importante. Endireitou os ombros, refez o rosto e entrou numa loja onde nunca fora bem atendida. Um casal abriu passagem. O vendedor interrompeu outra cliente para oferecer ajuda e chamou Helena de senhora antes que ela dissesse qualquer coisa. Ela pediu para experimentar o vestido mais caro apenas para vê-lo obedecer.

No provador, a luz branca revelou uma mancha na sua mão. Helena pensou que talvez tudo fosse apenas maquiagem da cabeça, e o encanto quebrou. O vestido continuava nela, mas o rosto antigo reapareceu no espelho. Quando saiu, o vendedor levou alguns segundos para reconhecê-la e perguntou, com menos gentileza, se ela pretendia comprar. Helena deixou a roupa no chão. Na calçada, precisou se apoiar numa banca: o poder cobrava fé ininterrupta, justamente o que ela passara a vida inteira sem ter em si mesma.

\begin{center}
* * *
\end{center}

À noite, o \textbf{garçom} fechou o La Fontaine.  
O salão vazio refletia seu cansaço.  
Sentou-se no balcão e serviu a si mesmo um resto de vinho barato.  
O reflexo na garrafa devolveu um rosto comprimido e ridículo. O garçom repetiu diante dele as perguntas que fizera durante o dia. “Tem certeza?” “Foi isso mesmo?” “Você não está confundindo?” Na terceira, começou a desconfiar de que sua descoberta também pudesse ser coincidência. A mão apertou o gargalo até os dedos doerem.

Sorriu, sem mostrar os dentes.
— Não — disse ao reflexo. — Eu sei o que fiz.

Não permitiria que a própria dúvida lhe roubasse aquilo.

\begin{center}
* * *
\end{center}

Do outro lado da cidade, \textbf{Helena} não conseguia dormir.  
O corpo ainda vibrava.  
A lembrança da transformação a fazia sorrir sozinha.  
No escuro, imaginou seu rosto iluminado, os cabelos loiros, a pele perfeita.  
No espelho do quarto, a mecha loira voltou. Helena não comemorou; respirou fundo e fez surgir a segunda, depois a terceira.

Ela se sentou na cama e, pela primeira vez, não se sentiu invisível. Pegou o telefone, abriu a agenda e escolheu um restaurante onde antes teria vergonha de entrar sozinha. Fez uma reserva para a noite seguinte usando apenas o primeiro nome. Tudo poderia ser real, desde que sustentasse a imagem até o fim.

\begin{center}
* * *
\end{center}

O \textbf{garçom} andava pelas ruas vazias.  
Cada passo fazia o ar esfriar.  
Passou por um grupo de jovens na calçada — riam alto, confiantes.  
Ele parou e apenas os olhou.  
A risada morreu.  
Um dos garotos baixou o olhar, outro fingiu atender o celular.  
Ele passou por eles e o silêncio ficou.  
Sorriu, satisfeito.

Não precisou encostar em ninguém nem levantar a voz. Bastou descobrir onde cada um guardava a própria insegurança e empurrar um pouco. Um dos jovens ainda tentou rir depois que ele passou, mas os outros perguntaram do que estava rindo. O garçom ouviu a discussão começar e guardou aquela reação como quem decora uma receita.

\begin{center}
* * *
\end{center}

Naquela madrugada, Helena ficou diante do espelho repetindo a aparência que usaria na reserva. Cada vez que o rosto falhava, recomeçava pelo formato dos olhos. No La Fontaine, o garçom anotou num bloco as perguntas que haviam funcionado e as reações que provocaram. Não conhecia o nome do que fazia, mas já planejava testar no dia seguinte se uma certeza antiga quebrava com a mesma facilidade que uma escolha de vinho.

Ela precisava acreditar sem interrupção. Ele precisava encontrar a interrupção na crença dos outros. Ainda não sabiam que essas duas descobertas acabariam na mesma mesa, mas ambos foram dormir desejando uma nova oportunidade de experimentar.
